\section{Discussion}

The central result is a separation between structural opportunity and structural
intervention effect. In the broad history, heterogeneous signs appear across
protocols and readouts, but the historical data cannot identify one common
opportunity treatment. In the narrower V21 counterfactual, topology-dependent
selection still changes sign after generic intervention and selection noise are
matched. The defensible conclusion is a conditional, sign-heterogeneous
selection effect: how and where structure intervenes matters at least as much
as whether a structural signal can be measured.

This result changes what a next method would need to demonstrate. A more
elaborate selector is not justified merely because a previous selector failed.
A convincing method would need to show, under a matched generic-intervention
control, that its policy changes the primary readout in a stable direction. It
would also need to separate label-free training evidence from post-hoc semantic
diagnostics. The current study does not propose that method; it supplies the
failure-localization protocol needed before designing one.

The findings also explain why local geometry, reconstruction difficulty, and
assignment stability should not be used interchangeably. A local metric can be
post-treatment and label-aware \citep{li2025dcboost}, a reconstruction loss can
reward a difficult coordinate without improving the partition
\citep{yan2026scmib}, and a training head can disagree with a clean-embedding
KMeans readout. Treating these as the same utility signal would recreate the
experimental loop this study is designed to stop \citep{li2025lfss}.
